\chapter{Glossary of Concepts}
\label{app:F}

The front-matter List of Abbreviations gives a quick lookup for every acronym used in this thesis. This appendix provides prose explanations of concepts that recur across multiple chapters.

\begin{description}
\item[Bus matrix] The internal interconnect fabric of the microcontroller through which the CPU core, DMA controller, and every memory-mapped peripheral reach the same underlying memory and register space. Section~3.2 describes the specific two-bus-master arrangement (CPU + DMA) that creates the cache coherency hazard addressed in CS02.

\item[Cache coherency] The property that every bus master observes a consistent view of memory. A CPU data cache alone is coherent; coherency becomes a live concern once a second, independent bus master (DMA controller) can modify memory the cache has already copied, as in CS02.

\item[Circular buffer mode] A DMA configuration in which the transfer address wraps around to the base after reaching the end of the buffer, enabling continuous reception without software reload. Used for USART1 MAVLink reception in CS02.

\item[DMB / ISB] Data Memory Barrier and Instruction Memory Barrier --- ARM instructions that enforce ordering of memory accesses and pipeline flushes. Critical for lock-free ring buffer correctness (CS04).

\item[EKF innovation] The difference between a sensor measurement and the measurement predicted by the current state estimate. Large innovations indicate model mismatch or outlier measurements. The innovation gate (CS05) rejects measurements with statistically improbable innovations.

\item[FPU] Floating-Point Unit --- the Cortex-M7 hardware floating-point accelerator (FPv5-D16). Must be enabled in Reset\_Handler via CPACR before any floating-point instruction executes; otherwise a UsageFault occurs (CS11).

\item[Input capture] A timer mode in which the current counter value is latched into a register when a GPIO edge is detected, providing hardware-precise timestamping without software latency. Used for radar frame-sync in CS03.

\item[MAVLink] Micro Air Vehicle Link --- a lightweight telemetry protocol for unmanned vehicles. This thesis uses MAVLink v2 messages: ATTITUDE (\#30), LOCAL\_POSITION\_NED (\#32), and GLOBAL\_POSITION\_INT (\#33).

\item[MPU] Memory Protection Unit --- configures per-region access and execute permissions. An XN (Execute-Never) region triggers a MemManage fault on instruction fetch (CS09).

\item[Quaternion] A 4-element representation of 3D rotation $[q_w, q_x, q_y, q_z]$ that avoids gimbal lock and composes efficiently. The conjugate $\mathbf{q}^{-1} = [q_w, -q_x, -q_y, -q_z]$ represents the inverse rotation; applying it when the forward rotation is needed (or vice versa) produces incorrect results (CS06).

\item[Renode] Open-source functional simulator by Antmicro. Used as the primary verification platform. Diverges from physical silicon in specific, documented ways (Chapter~5).

\item[Struct packing] Compiler attribute that eliminates padding between struct members, required when the struct layout must match an external wire protocol byte-for-byte (CS02).
\end{description}